\documentclass[10pt,a4paper]{article}
%% ВЕРСИЯ: номер и история — в SESSIYA.md арки zhurnal/2026-09-27_listok-18-veroyatnost; собрано sobrat.py из alef-istochnik/telo.tex + kartinki.py
%% План — голосовое владельца 30.09 18:09 (zhurnal/2026-09-27_listok-18-veroyatnost/GOLOS-2026-09-30-1809-plan-alef.md).
%% Ответы проверены перебором: listok-veroyatnost/prov_alef_v1.py (ALL OK). Вёрстка — как 18alpha-veroyatnost.tex.
\usepackage{iftex}
\ifPDFTeX
  \usepackage[T2A]{fontenc}
  \usepackage[utf8]{inputenc}
  \usepackage[russian]{babel}
\else
\usepackage{fontspec}
\usepackage{polyglossia}
\setmainlanguage{russian}
\setotherlanguage{english}
\setmainfont{SourceSerifPro}[
  Path           = ../listok-derevya/shrifty/,
  Extension      = .ttf,
  UprightFont    = *_400Regular,
  ItalicFont     = *_400Regular_Italic,
  BoldFont       = *_700Bold,
  BoldItalicFont = *_700Bold_Italic]
\fi
\usepackage[left=11mm,right=13mm,top=18mm,bottom=14mm]{geometry}
\usepackage{amsmath,amssymb}
\usepackage{tikz}
\usepackage{wrapfig}
\ifLuaTeX
  \usepackage[protrusion=true,expansion=true,tracking=true,
              stretch=50,shrink=30,final]{microtype}
\else
  \usepackage[protrusion=true,final]{microtype}
\fi
\emergencystretch=1em
\linespread{0.98}
\frenchspacing
\setlength{\parindent}{0pt}
\pagestyle{empty}

\newlength{\shirinaokna}\setlength{\shirinaokna}{12mm}
\newlength{\vysotaokna}\setlength{\vysotaokna}{0.7\baselineskip}
\newcommand{\okno}{\framebox[\shirinaokna]{\rule{0pt}{\vysotaokna}}}
\newlength{\shirinanomera}\setlength{\shirinanomera}{9mm}
\newlength{\zazor}\setlength{\zazor}{0.6em}
\newcommand{\blok}[1]{\par\noindent #1\par}
\newcommand{\nomer}[2]{\makebox[\shirinanomera][l]{%
  \fbox{\rule{0pt}{\vysotaokna}\textbf{#1}$^{#2}$}}}
\newcommand{\term}[1]{\textit{#1}}
\newcommand{\zvezdy}{%
  \par\medskip
  \centerline{\rule{0.3\linewidth}{0.4pt}\quad $*$\ \ $*$\ \ $*$\quad
              \rule{0.3\linewidth}{0.4pt}}%
  \par\nopagebreak\medskip}
\newcommand{\zad}[3]{\begin{samepage}\blok{\nomer{#1}{#2}\hspace{\zazor}\okno\hspace{\zazor}#3}\end{samepage}\par\vspace{9pt}}
%% \zadp{номер}{маркер} — задача без общего условия: номер, и на той же строке сразу первый \punkt. perevod.py понимает.
\newif\ifvstroku
\newcommand{\zadp}[2]{\par\noindent\nomer{#1}{#2}\hspace{\zazor}\global\vstrokutrue}
\newcommand{\punkt}[3]{\ifvstroku\global\vstrokufalse\else\par\nopagebreak\noindent\fi\okno\hspace{\zazor}\textbf{#1)}$^{#2}$\ #3\par}
%% Текст и картинка рядом: \sboku{ширина картинки}{картинка}{текст задач}. perevod.py разворачивает это для коллег в «текст, затем картинка по центру».
\newcommand{\sboku}[3]{\par\noindent\begin{minipage}[t]{\dimexpr\linewidth-#1-1.2em\relax}\vspace{0pt}#3\end{minipage}\hfill\begin{minipage}[t]{#1}\vspace{0pt}\centering #2\end{minipage}\par\vspace{9pt}}
%% \celoe{…} — неразрывный блок (задача вместе с картинкой переносится целиком); \risunok{…} — картинка по центру с отступами;
%% \obtek[строк]{ширина}{картинка} — обтекание справа (wrapfig), ставится перед абзацем. perevod.py понимает все три.
\newcommand{\celoe}[1]{\par\noindent\begin{minipage}{\linewidth}#1\end{minipage}\par\vspace{9pt}}
\newcommand{\risunok}[1]{\par\nopagebreak\medskip\centerline{#1}\par\smallskip}
\newcommand{\obtek}[3][]{\begin{wrapfigure}[#1]{r}{#2}\centering #3\end{wrapfigure}}
\newcommand{\tekst}[1]{\blok{#1}\par\vspace{9pt}}
%% Поле для подписи на третьей странице
\AddToHook{shipout/foreground}{\ifnum\value{page}=3 \put(305,-22){\small Фамилия, имя:\ \rule{6.2cm}{0.4pt}}\fi}
%% Номер страницы внизу по центру
\AddToHook{shipout/foreground}{\put(298.8,-818){\makebox[0pt]{\small\thepage}}}
\begin{document}
\noindent\rlap{\makebox[\linewidth]{\textit{Школа № 179. Ключики 2026/2027 уч.~год}}}\hfill{\LARGE$\aleph$}\par\medskip
\centerline{\rule{0.62\linewidth}{0.4pt}}\par\smallskip
\centerline{\large\textbf{18. \textsc{Вероятность}}}\par\smallskip

\zad{1}{}{На выборах за кандидата А подано $a$ голосов, за Б~— $b$ голосов, $a>b$. Бюллетени перемешали и достают по одному, пока не достанут все, и каждый раз объявляют счёт. С какой вероятностью А после каждого бюллетеня был строго впереди? \par\vspace{3pt}К этой задаче мы ещё не раз вернёмся.}

\begin{samepage}
\zadp{2}{\circ}
\punkt{а}{}{Какие здесь элементарные исходы? Сколько их и чему равна вероятность каждого?}
\punkt{б}{}{Найдите ответ в задаче 1 при $a=3$, $b=2$ и при $a=4$, $b=2$.}
\punkt{в}{}{Как изобразить исход, чтобы с первого взгляда было видно, что А всё время был впереди?}
\end{samepage}\par\vspace{9pt}

\celoe{\tekst{\term{Путь}~— ломаная из звеньев $\nearrow$ и $\searrow$, \term{длина} пути~— число звеньев. Если начало пути не указано, он начинается в $(0,0)$. Вот десять путей из $(0,0)$ в $(8,2)$. Сколько всего таких путей?}\risunok{\begin{tikzpicture}[scale=0.3] \draw[step=1,gray!28,line width=0.3pt] (0,-3) grid (8,5); \draw[gray!70,line width=0.6pt] (0,0) -- (8,0); \draw[line width=1.3pt,line join=round,line cap=round] (0,0) -- (1,1) -- (2,2) -- (3,3) -- (4,4) -- (5,3) -- (6,4) -- (7,3) -- (8,2); \draw[step=1,gray!28,line width=0.3pt] (10.6,-3) grid (18.6,5); \draw[gray!70,line width=0.6pt] (10.6,0) -- (18.6,0); \draw[line width=1.3pt,line join=round,line cap=round] (10.6,0) -- (11.6,-1) -- (12.6,0) -- (13.6,1) -- (14.6,0) -- (15.6,1) -- (16.6,0) -- (17.6,1) -- (18.6,2); \draw[step=1,gray!28,line width=0.3pt] (21.2,-3) grid (29.2,5); \draw[gray!70,line width=0.6pt] (21.2,0) -- (29.2,0); \draw[line width=1.3pt,line join=round,line cap=round] (21.2,0) -- (22.2,-1) -- (23.2,0) -- (24.2,-1) -- (25.2,0) -- (26.2,-1) -- (27.2,0) -- (28.2,1) -- (29.2,2); \draw[step=1,gray!28,line width=0.3pt] (31.8,-3) grid (39.8,5); \draw[gray!70,line width=0.6pt] (31.8,0) -- (39.8,0); \draw[line width=1.3pt,line join=round,line cap=round] (31.8,0) -- (32.8,1) -- (33.8,2) -- (34.8,3) -- (35.8,4) -- (36.8,3) -- (37.8,2) -- (38.8,1) -- (39.8,2); \draw[step=1,gray!28,line width=0.3pt] (42.4,-3) grid (50.4,5); \draw[gray!70,line width=0.6pt] (42.4,0) -- (50.4,0); \draw[line width=1.3pt,line join=round,line cap=round] (42.4,0) -- (43.4,-1) -- (44.4,-2) -- (45.4,-1) -- (46.4,0) -- (47.4,1) -- (48.4,2) -- (49.4,1) -- (50.4,2); \draw[step=1,gray!28,line width=0.3pt] (0,-12.2) grid (8,-4.2); \draw[gray!70,line width=0.6pt] (0,-9.2) -- (8,-9.2); \draw[line width=1.3pt,line join=round,line cap=round] (0,-9.2) -- (1,-10.2) -- (2,-11.2) -- (3,-10.2) -- (4,-9.2) -- (5,-8.2) -- (6,-9.2) -- (7,-8.2) -- (8,-7.2); \draw[step=1,gray!28,line width=0.3pt] (10.6,-12.2) grid (18.6,-4.2); \draw[gray!70,line width=0.6pt] (10.6,-9.2) -- (18.6,-9.2); \draw[line width=1.3pt,line join=round,line cap=round] (10.6,-9.2) -- (11.6,-8.2) -- (12.6,-9.2) -- (13.6,-8.2) -- (14.6,-7.2) -- (15.6,-6.2) -- (16.6,-7.2) -- (17.6,-8.2) -- (18.6,-7.2); \draw[step=1,gray!28,line width=0.3pt] (21.2,-12.2) grid (29.2,-4.2); \draw[gray!70,line width=0.6pt] (21.2,-9.2) -- (29.2,-9.2); \draw[line width=1.3pt,line join=round,line cap=round] (21.2,-9.2) -- (22.2,-10.2) -- (23.2,-9.2) -- (24.2,-8.2) -- (25.2,-7.2) -- (26.2,-6.2) -- (27.2,-7.2) -- (28.2,-6.2) -- (29.2,-7.2); \draw[step=1,gray!28,line width=0.3pt] (31.8,-12.2) grid (39.8,-4.2); \draw[gray!70,line width=0.6pt] (31.8,-9.2) -- (39.8,-9.2); \draw[line width=1.3pt,line join=round,line cap=round] (31.8,-9.2) -- (32.8,-8.2) -- (33.8,-9.2) -- (34.8,-8.2) -- (35.8,-9.2) -- (36.8,-8.2) -- (37.8,-7.2) -- (38.8,-8.2) -- (39.8,-7.2); \draw[step=1,gray!28,line width=0.3pt] (42.4,-12.2) grid (50.4,-4.2); \draw[gray!70,line width=0.6pt] (42.4,-9.2) -- (50.4,-9.2); \draw[line width=1.3pt,line join=round,line cap=round] (42.4,-9.2) -- (43.4,-8.2) -- (44.4,-7.2) -- (45.4,-6.2) -- (46.4,-7.2) -- (47.4,-6.2) -- (48.4,-7.2) -- (49.4,-6.2) -- (50.4,-7.2); \end{tikzpicture}}}

\risunok{\begin{tikzpicture}[xscale=1.4,yscale=0.48,>=stealth] \draw[gray!70,line width=0.6pt,->] (0,0) -- (12.6,0); \draw[gray!70,line width=0.6pt,->] (0,0) -- (0,12.9); \draw[->,black!55,line width=0.5pt,shorten >=8.6pt,shorten <=8.6pt] (0,0) -- (1,1); \draw[->,black!55,line width=0.5pt,shorten >=8.6pt,shorten <=8.6pt] (1,1) -- (2,2); \draw[->,black!55,line width=0.5pt,shorten >=8.6pt,shorten <=8.6pt] (1,1) -- (2,0); \draw[->,black!55,line width=0.5pt,shorten >=8.6pt,shorten <=8.6pt] (2,0) -- (3,1); \draw[->,black!55,line width=0.5pt,shorten >=8.6pt,shorten <=8.6pt] (2,2) -- (3,3); \draw[->,black!55,line width=0.5pt,shorten >=8.6pt,shorten <=8.6pt] (2,2) -- (3,1); \draw[->,black!55,line width=0.5pt,shorten >=8.6pt,shorten <=8.6pt] (3,1) -- (4,2); \draw[->,black!55,line width=0.5pt,shorten >=8.6pt,shorten <=8.6pt] (3,1) -- (4,0); \draw[->,black!55,line width=0.5pt,shorten >=8.6pt,shorten <=8.6pt] (3,3) -- (4,4); \draw[->,black!55,line width=0.5pt,shorten >=8.6pt,shorten <=8.6pt] (3,3) -- (4,2); \draw[->,black!55,line width=0.5pt,shorten >=8.6pt,shorten <=8.6pt] (4,0) -- (5,1); \draw[->,black!55,line width=0.5pt,shorten >=8.6pt,shorten <=8.6pt] (4,2) -- (5,3); \draw[->,black!55,line width=0.5pt,shorten >=8.6pt,shorten <=8.6pt] (4,2) -- (5,1); \draw[->,black!55,line width=0.5pt,shorten >=8.6pt,shorten <=8.6pt] (4,4) -- (5,5); \draw[->,black!55,line width=0.5pt,shorten >=8.6pt,shorten <=8.6pt] (4,4) -- (5,3); \draw[->,black!55,line width=0.5pt,shorten >=8.6pt,shorten <=8.6pt] (5,1) -- (6,2); \draw[->,black!55,line width=0.5pt,shorten >=8.6pt,shorten <=8.6pt] (5,1) -- (6,0); \draw[->,black!55,line width=0.5pt,shorten >=8.6pt,shorten <=8.6pt] (5,3) -- (6,4); \draw[->,black!55,line width=0.5pt,shorten >=8.6pt,shorten <=8.6pt] (5,3) -- (6,2); \draw[->,black!55,line width=0.5pt,shorten >=8.6pt,shorten <=8.6pt] (5,5) -- (6,6); \draw[->,black!55,line width=0.5pt,shorten >=8.6pt,shorten <=8.6pt] (5,5) -- (6,4); \draw[->,black!55,line width=0.5pt,shorten >=8.6pt,shorten <=8.6pt] (6,0) -- (7,1); \draw[->,black!55,line width=0.5pt,shorten >=8.6pt,shorten <=8.6pt] (6,2) -- (7,3); \draw[->,black!55,line width=0.5pt,shorten >=8.6pt,shorten <=8.6pt] (6,2) -- (7,1); \draw[->,black!55,line width=0.5pt,shorten >=8.6pt,shorten <=8.6pt] (6,4) -- (7,5); \draw[->,black!55,line width=0.5pt,shorten >=8.6pt,shorten <=8.6pt] (6,4) -- (7,3); \draw[->,black!55,line width=0.5pt,shorten >=8.6pt,shorten <=8.6pt] (6,6) -- (7,7); \draw[->,black!55,line width=0.5pt,shorten >=8.6pt,shorten <=8.6pt] (6,6) -- (7,5); \draw[->,black!55,line width=0.5pt,shorten >=8.6pt,shorten <=8.6pt] (7,1) -- (8,2); \draw[->,black!55,line width=0.5pt,shorten >=8.6pt,shorten <=8.6pt] (7,1) -- (8,0); \draw[->,black!55,line width=0.5pt,shorten >=8.6pt,shorten <=8.6pt] (7,3) -- (8,4); \draw[->,black!55,line width=0.5pt,shorten >=8.6pt,shorten <=8.6pt] (7,3) -- (8,2); \draw[->,black!55,line width=0.5pt,shorten >=8.6pt,shorten <=8.6pt] (7,5) -- (8,6); \draw[->,black!55,line width=0.5pt,shorten >=8.6pt,shorten <=8.6pt] (7,5) -- (8,4); \draw[->,black!55,line width=0.5pt,shorten >=8.6pt,shorten <=8.6pt] (7,7) -- (8,8); \draw[->,black!55,line width=0.5pt,shorten >=8.6pt,shorten <=8.6pt] (7,7) -- (8,6); \draw[->,black!55,line width=0.5pt,shorten >=8.6pt,shorten <=8.6pt] (8,0) -- (9,1); \draw[->,black!55,line width=0.5pt,shorten >=8.6pt,shorten <=8.6pt] (8,2) -- (9,3); \draw[->,black!55,line width=0.5pt,shorten >=8.6pt,shorten <=8.6pt] (8,2) -- (9,1); \draw[->,black!55,line width=0.5pt,shorten >=8.6pt,shorten <=8.6pt] (8,4) -- (9,5); \draw[->,black!55,line width=0.5pt,shorten >=8.6pt,shorten <=8.6pt] (8,4) -- (9,3); \draw[->,black!55,line width=0.5pt,shorten >=8.6pt,shorten <=8.6pt] (8,6) -- (9,7); \draw[->,black!55,line width=0.5pt,shorten >=8.6pt,shorten <=8.6pt] (8,6) -- (9,5); \draw[->,black!55,line width=0.5pt,shorten >=8.6pt,shorten <=8.6pt] (8,8) -- (9,9); \draw[->,black!55,line width=0.5pt,shorten >=8.6pt,shorten <=8.6pt] (8,8) -- (9,7); \draw[->,black!55,line width=0.5pt,shorten >=8.6pt,shorten <=8.6pt] (9,1) -- (10,2); \draw[->,black!55,line width=0.5pt,shorten >=8.6pt,shorten <=8.6pt] (9,1) -- (10,0); \draw[->,black!55,line width=0.5pt,shorten >=8.6pt,shorten <=8.6pt] (9,3) -- (10,4); \draw[->,black!55,line width=0.5pt,shorten >=8.6pt,shorten <=8.6pt] (9,3) -- (10,2); \draw[->,black!55,line width=0.5pt,shorten >=8.6pt,shorten <=8.6pt] (9,5) -- (10,6); \draw[->,black!55,line width=0.5pt,shorten >=8.6pt,shorten <=8.6pt] (9,5) -- (10,4); \draw[->,black!55,line width=0.5pt,shorten >=8.6pt,shorten <=8.6pt] (9,7) -- (10,8); \draw[->,black!55,line width=0.5pt,shorten >=8.6pt,shorten <=8.6pt] (9,7) -- (10,6); \draw[->,black!55,line width=0.5pt,shorten >=8.6pt,shorten <=8.6pt] (9,9) -- (10,10); \draw[->,black!55,line width=0.5pt,shorten >=8.6pt,shorten <=8.6pt] (9,9) -- (10,8); \draw[->,black!55,line width=0.5pt,shorten >=8.6pt,shorten <=8.6pt] (10,0) -- (11,1); \draw[->,black!55,line width=0.5pt,shorten >=8.6pt,shorten <=8.6pt] (10,2) -- (11,3); \draw[->,black!55,line width=0.5pt,shorten >=8.6pt,shorten <=8.6pt] (10,2) -- (11,1); \draw[->,black!55,line width=0.5pt,shorten >=8.6pt,shorten <=8.6pt] (10,4) -- (11,5); \draw[->,black!55,line width=0.5pt,shorten >=8.6pt,shorten <=8.6pt] (10,4) -- (11,3); \draw[->,black!55,line width=0.5pt,shorten >=8.6pt,shorten <=8.6pt] (10,6) -- (11,7); \draw[->,black!55,line width=0.5pt,shorten >=8.6pt,shorten <=8.6pt] (10,6) -- (11,5); \draw[->,black!55,line width=0.5pt,shorten >=8.6pt,shorten <=8.6pt] (10,8) -- (11,9); \draw[->,black!55,line width=0.5pt,shorten >=8.6pt,shorten <=8.6pt] (10,8) -- (11,7); \draw[->,black!55,line width=0.5pt,shorten >=8.6pt,shorten <=8.6pt] (10,10) -- (11,11); \draw[->,black!55,line width=0.5pt,shorten >=8.6pt,shorten <=8.6pt] (10,10) -- (11,9); \draw[->,black!55,line width=0.5pt,shorten >=8.6pt,shorten <=8.6pt] (11,1) -- (12,2); \draw[->,black!55,line width=0.5pt,shorten >=8.6pt,shorten <=8.6pt] (11,1) -- (12,0); \draw[->,black!55,line width=0.5pt,shorten >=8.6pt,shorten <=8.6pt] (11,3) -- (12,4); \draw[->,black!55,line width=0.5pt,shorten >=8.6pt,shorten <=8.6pt] (11,3) -- (12,2); \draw[->,black!55,line width=0.5pt,shorten >=8.6pt,shorten <=8.6pt] (11,5) -- (12,6); \draw[->,black!55,line width=0.5pt,shorten >=8.6pt,shorten <=8.6pt] (11,5) -- (12,4); \draw[->,black!55,line width=0.5pt,shorten >=8.6pt,shorten <=8.6pt] (11,7) -- (12,8); \draw[->,black!55,line width=0.5pt,shorten >=8.6pt,shorten <=8.6pt] (11,7) -- (12,6); \draw[->,black!55,line width=0.5pt,shorten >=8.6pt,shorten <=8.6pt] (11,9) -- (12,10); \draw[->,black!55,line width=0.5pt,shorten >=8.6pt,shorten <=8.6pt] (11,9) -- (12,8); \draw[->,black!55,line width=0.5pt,shorten >=8.6pt,shorten <=8.6pt] (11,11) -- (12,12); \draw[->,black!55,line width=0.5pt,shorten >=8.6pt,shorten <=8.6pt] (11,11) -- (12,10); \node[circle,draw,line width=0.45pt,inner sep=0pt,minimum size=16pt,font=\footnotesize,fill=white] at (0,0) {1}; \node[circle,draw,line width=0.45pt,inner sep=0pt,minimum size=16pt,font=\footnotesize,fill=white] at (1,1) {1}; \node[circle,draw,line width=0.45pt,inner sep=0pt,minimum size=16pt,font=\footnotesize,fill=white] at (2,0) {1}; \node[circle,draw,line width=0.45pt,inner sep=0pt,minimum size=16pt,font=\footnotesize,fill=white] at (2,2) {1}; \node[circle,draw,line width=0.45pt,inner sep=0pt,minimum size=16pt,font=\footnotesize,fill=white] at (3,1) {2}; \node[circle,draw,line width=0.45pt,inner sep=0pt,minimum size=16pt,font=\footnotesize,fill=white] at (3,3) {1}; \node[circle,draw,line width=0.45pt,inner sep=0pt,minimum size=16pt,font=\footnotesize,fill=white] at (4,0) {2}; \node[circle,draw,line width=0.45pt,inner sep=0pt,minimum size=16pt,font=\footnotesize,fill=white] at (4,2) {3}; \node[circle,draw,line width=0.45pt,inner sep=0pt,minimum size=16pt,font=\footnotesize,fill=white] at (4,4) {1}; \node[circle,draw,line width=0.45pt,inner sep=0pt,minimum size=16pt,font=\footnotesize,fill=white] at (5,1) {5}; \node[circle,draw,line width=0.45pt,inner sep=0pt,minimum size=16pt,font=\footnotesize,fill=white] at (5,3) {4}; \node[circle,draw,line width=0.45pt,inner sep=0pt,minimum size=16pt,font=\footnotesize,fill=white] at (5,5) {1}; \node[circle,draw,line width=0.45pt,inner sep=0pt,minimum size=16pt,font=\footnotesize,fill=white] at (6,0) {}; \node[circle,draw,line width=0.45pt,inner sep=0pt,minimum size=16pt,font=\footnotesize,fill=white] at (6,2) {}; \node[circle,draw,line width=0.45pt,inner sep=0pt,minimum size=16pt,font=\footnotesize,fill=white] at (6,4) {}; \node[circle,draw,line width=0.45pt,inner sep=0pt,minimum size=16pt,font=\footnotesize,fill=white] at (6,6) {}; \node[circle,draw,line width=0.45pt,inner sep=0pt,minimum size=16pt,font=\footnotesize,fill=white] at (7,1) {}; \node[circle,draw,line width=0.45pt,inner sep=0pt,minimum size=16pt,font=\footnotesize,fill=white] at (7,3) {}; \node[circle,draw,line width=0.45pt,inner sep=0pt,minimum size=16pt,font=\footnotesize,fill=white] at (7,5) {}; \node[circle,draw,line width=0.45pt,inner sep=0pt,minimum size=16pt,font=\footnotesize,fill=white] at (7,7) {}; \node[circle,draw,line width=0.45pt,inner sep=0pt,minimum size=16pt,font=\footnotesize,fill=white] at (8,0) {}; \node[circle,draw,line width=0.45pt,inner sep=0pt,minimum size=16pt,font=\footnotesize,fill=white] at (8,2) {}; \node[circle,draw,line width=0.45pt,inner sep=0pt,minimum size=16pt,font=\footnotesize,fill=white] at (8,4) {}; \node[circle,draw,line width=0.45pt,inner sep=0pt,minimum size=16pt,font=\footnotesize,fill=white] at (8,6) {}; \node[circle,draw,line width=0.45pt,inner sep=0pt,minimum size=16pt,font=\footnotesize,fill=white] at (8,8) {}; \node[circle,draw,line width=0.45pt,inner sep=0pt,minimum size=16pt,font=\footnotesize,fill=white] at (9,1) {}; \node[circle,draw,line width=0.45pt,inner sep=0pt,minimum size=16pt,font=\footnotesize,fill=white] at (9,3) {}; \node[circle,draw,line width=0.45pt,inner sep=0pt,minimum size=16pt,font=\footnotesize,fill=white] at (9,5) {}; \node[circle,draw,line width=0.45pt,inner sep=0pt,minimum size=16pt,font=\footnotesize,fill=white] at (9,7) {}; \node[circle,draw,line width=0.45pt,inner sep=0pt,minimum size=16pt,font=\footnotesize,fill=white] at (9,9) {}; \node[circle,draw,line width=0.45pt,inner sep=0pt,minimum size=16pt,font=\footnotesize,fill=white] at (10,0) {}; \node[circle,draw,line width=0.45pt,inner sep=0pt,minimum size=16pt,font=\footnotesize,fill=white] at (10,2) {}; \node[circle,draw,line width=0.45pt,inner sep=0pt,minimum size=16pt,font=\footnotesize,fill=white] at (10,4) {}; \node[circle,draw,line width=0.45pt,inner sep=0pt,minimum size=16pt,font=\footnotesize,fill=white] at (10,6) {}; \node[circle,draw,line width=0.45pt,inner sep=0pt,minimum size=16pt,font=\footnotesize,fill=white] at (10,8) {}; \node[circle,draw,line width=0.45pt,inner sep=0pt,minimum size=16pt,font=\footnotesize,fill=white] at (10,10) {}; \node[circle,draw,line width=0.45pt,inner sep=0pt,minimum size=16pt,font=\footnotesize,fill=white] at (11,1) {}; \node[circle,draw,line width=0.45pt,inner sep=0pt,minimum size=16pt,font=\footnotesize,fill=white] at (11,3) {}; \node[circle,draw,line width=0.45pt,inner sep=0pt,minimum size=16pt,font=\footnotesize,fill=white] at (11,5) {}; \node[circle,draw,line width=0.45pt,inner sep=0pt,minimum size=16pt,font=\footnotesize,fill=white] at (11,7) {}; \node[circle,draw,line width=0.45pt,inner sep=0pt,minimum size=16pt,font=\footnotesize,fill=white] at (11,9) {}; \node[circle,draw,line width=0.45pt,inner sep=0pt,minimum size=16pt,font=\footnotesize,fill=white] at (11,11) {}; \node[circle,draw,line width=0.45pt,inner sep=0pt,minimum size=16pt,font=\footnotesize,fill=white] at (12,0) {}; \node[circle,draw,line width=0.45pt,inner sep=0pt,minimum size=16pt,font=\footnotesize,fill=white] at (12,2) {}; \node[circle,draw,line width=0.45pt,inner sep=0pt,minimum size=16pt,font=\footnotesize,fill=white] at (12,4) {}; \node[circle,draw,line width=0.45pt,inner sep=0pt,minimum size=16pt,font=\footnotesize,fill=white] at (12,6) {}; \node[circle,draw,line width=0.45pt,inner sep=0pt,minimum size=16pt,font=\footnotesize,fill=white] at (12,8) {}; \node[circle,draw,line width=0.45pt,inner sep=0pt,minimum size=16pt,font=\footnotesize,fill=white] at (12,10) {}; \node[circle,draw,line width=0.45pt,inner sep=0pt,minimum size=16pt,font=\footnotesize,fill=white] at (12,12) {}; \node[font=\scriptsize,gray] at (2,-0.95) {2}; \node[font=\scriptsize,gray] at (4,-0.95) {4}; \node[font=\scriptsize,gray] at (6,-0.95) {6}; \node[font=\scriptsize,gray] at (8,-0.95) {8}; \node[font=\scriptsize,gray] at (10,-0.95) {10}; \node[font=\scriptsize,gray] at (12,-0.95) {12}; \node[font=\scriptsize,gray,anchor=east] at (-0.12,2) {2}; \node[font=\scriptsize,gray,anchor=east] at (-0.12,4) {4}; \node[font=\scriptsize,gray,anchor=east] at (-0.12,6) {6}; \node[font=\scriptsize,gray,anchor=east] at (-0.12,8) {8}; \node[font=\scriptsize,gray,anchor=east] at (-0.12,10) {10}; \node[font=\scriptsize,gray,anchor=east] at (-0.12,12) {12}; \node[font=\scriptsize,gray,anchor=north east] at (-0.05,-0.25) {0}; \node[anchor=north west,inner sep=0pt] at (0.25,12.75) {\parbox[t]{15cm}{\parshape 9 0cm 15.00cm 0cm 14.15cm 0cm 12.94cm 0cm 11.73cm 0cm 10.52cm 0cm 9.31cm 0cm 8.10cm 0cm 6.89cm 0cm 5.68cm\noindent \textit{Путь Дика}~— путь из $(0,0)$, который ни разу не опускается ниже оси.\\Пути Дика длины не больше 12~— это пути по стрелкам схемы ниже.\\\textit{Треугольником Каталана} называют схему, в каждой точке которой записано число путей Дика, ведущих в эту точку. Часть чисел уже~вписана.}}; \end{tikzpicture}}\vspace{4pt}

\begin{samepage}
\zadp{3}{\circ}
\punkt{а}{}{В каких точках может кончаться путь Дика?}
\punkt{б}{}{Заполните треугольник Каталана. По какому правилу число в точке получается из соседних?}
\punkt{в}{}{Найдите ответ в задаче 1 при всех $a>b$, $a+b\leqslant 8$. Предположите, чему он равен в~общем~случае.}
\end{samepage}\par\vspace{9pt}

\celoe{\tekst{Число путей Дика в $(2n,0)$ называют \term{числом Каталана} и обозначают $C_n$. На рисунке~— все пути Дика~в~$(8,0)$.}\risunok{\begin{tikzpicture}[scale=0.27] \draw[step=1,gray!28,line width=0.3pt] (0,0) grid (8,4); \draw[gray!70,line width=0.6pt] (0,0) -- (8,0); \draw[line width=1.3pt,line join=round,line cap=round] (0,0) -- (1,1) -- (2,2) -- (3,3) -- (4,4) -- (5,3) -- (6,2) -- (7,1) -- (8,0); \draw[step=1,gray!28,line width=0.3pt] (9.6,0) grid (17.6,4); \draw[gray!70,line width=0.6pt] (9.6,0) -- (17.6,0); \draw[line width=1.3pt,line join=round,line cap=round] (9.6,0) -- (10.6,1) -- (11.6,2) -- (12.6,3) -- (13.6,2) -- (14.6,3) -- (15.6,2) -- (16.6,1) -- (17.6,0); \draw[step=1,gray!28,line width=0.3pt] (19.2,0) grid (27.2,4); \draw[gray!70,line width=0.6pt] (19.2,0) -- (27.2,0); \draw[line width=1.3pt,line join=round,line cap=round] (19.2,0) -- (20.2,1) -- (21.2,2) -- (22.2,3) -- (23.2,2) -- (24.2,1) -- (25.2,2) -- (26.2,1) -- (27.2,0); \draw[step=1,gray!28,line width=0.3pt] (28.8,0) grid (36.8,4); \draw[gray!70,line width=0.6pt] (28.8,0) -- (36.8,0); \draw[line width=1.3pt,line join=round,line cap=round] (28.8,0) -- (29.8,1) -- (30.8,2) -- (31.8,3) -- (32.8,2) -- (33.8,1) -- (34.8,0) -- (35.8,1) -- (36.8,0); \draw[step=1,gray!28,line width=0.3pt] (38.4,0) grid (46.4,4); \draw[gray!70,line width=0.6pt] (38.4,0) -- (46.4,0); \draw[line width=1.3pt,line join=round,line cap=round] (38.4,0) -- (39.4,1) -- (40.4,2) -- (41.4,1) -- (42.4,2) -- (43.4,3) -- (44.4,2) -- (45.4,1) -- (46.4,0); \draw[step=1,gray!28,line width=0.3pt] (48,0) grid (56,4); \draw[gray!70,line width=0.6pt] (48,0) -- (56,0); \draw[line width=1.3pt,line join=round,line cap=round] (48,0) -- (49,1) -- (50,2) -- (51,1) -- (52,2) -- (53,1) -- (54,2) -- (55,1) -- (56,0); \draw[step=1,gray!28,line width=0.3pt] (57.6,0) grid (65.6,4); \draw[gray!70,line width=0.6pt] (57.6,0) -- (65.6,0); \draw[line width=1.3pt,line join=round,line cap=round] (57.6,0) -- (58.6,1) -- (59.6,2) -- (60.6,1) -- (61.6,2) -- (62.6,1) -- (63.6,0) -- (64.6,1) -- (65.6,0); \draw[step=1,gray!28,line width=0.3pt] (0,-5.8) grid (8,-1.8); \draw[gray!70,line width=0.6pt] (0,-5.8) -- (8,-5.8); \draw[line width=1.3pt,line join=round,line cap=round] (0,-5.8) -- (1,-4.8) -- (2,-3.8) -- (3,-4.8) -- (4,-5.8) -- (5,-4.8) -- (6,-3.8) -- (7,-4.8) -- (8,-5.8); \draw[step=1,gray!28,line width=0.3pt] (9.6,-5.8) grid (17.6,-1.8); \draw[gray!70,line width=0.6pt] (9.6,-5.8) -- (17.6,-5.8); \draw[line width=1.3pt,line join=round,line cap=round] (9.6,-5.8) -- (10.6,-4.8) -- (11.6,-3.8) -- (12.6,-4.8) -- (13.6,-5.8) -- (14.6,-4.8) -- (15.6,-5.8) -- (16.6,-4.8) -- (17.6,-5.8); \draw[step=1,gray!28,line width=0.3pt] (19.2,-5.8) grid (27.2,-1.8); \draw[gray!70,line width=0.6pt] (19.2,-5.8) -- (27.2,-5.8); \draw[line width=1.3pt,line join=round,line cap=round] (19.2,-5.8) -- (20.2,-4.8) -- (21.2,-5.8) -- (22.2,-4.8) -- (23.2,-3.8) -- (24.2,-2.8) -- (25.2,-3.8) -- (26.2,-4.8) -- (27.2,-5.8); \draw[step=1,gray!28,line width=0.3pt] (28.8,-5.8) grid (36.8,-1.8); \draw[gray!70,line width=0.6pt] (28.8,-5.8) -- (36.8,-5.8); \draw[line width=1.3pt,line join=round,line cap=round] (28.8,-5.8) -- (29.8,-4.8) -- (30.8,-5.8) -- (31.8,-4.8) -- (32.8,-3.8) -- (33.8,-4.8) -- (34.8,-3.8) -- (35.8,-4.8) -- (36.8,-5.8); \draw[step=1,gray!28,line width=0.3pt] (38.4,-5.8) grid (46.4,-1.8); \draw[gray!70,line width=0.6pt] (38.4,-5.8) -- (46.4,-5.8); \draw[line width=1.3pt,line join=round,line cap=round] (38.4,-5.8) -- (39.4,-4.8) -- (40.4,-5.8) -- (41.4,-4.8) -- (42.4,-3.8) -- (43.4,-4.8) -- (44.4,-5.8) -- (45.4,-4.8) -- (46.4,-5.8); \draw[step=1,gray!28,line width=0.3pt] (48,-5.8) grid (56,-1.8); \draw[gray!70,line width=0.6pt] (48,-5.8) -- (56,-5.8); \draw[line width=1.3pt,line join=round,line cap=round] (48,-5.8) -- (49,-4.8) -- (50,-5.8) -- (51,-4.8) -- (52,-5.8) -- (53,-4.8) -- (54,-3.8) -- (55,-4.8) -- (56,-5.8); \draw[step=1,gray!28,line width=0.3pt] (57.6,-5.8) grid (65.6,-1.8); \draw[gray!70,line width=0.6pt] (57.6,-5.8) -- (65.6,-5.8); \draw[line width=1.3pt,line join=round,line cap=round] (57.6,-5.8) -- (58.6,-4.8) -- (59.6,-5.8) -- (60.6,-4.8) -- (61.6,-5.8) -- (62.6,-4.8) -- (63.6,-5.8) -- (64.6,-4.8) -- (65.6,-5.8); \end{tikzpicture}}}

\begin{samepage}
\zadp{4}{\circ}
\punkt{а}{}{Скажем, что путь Дика впервые возвращается на ось в точке $(2k,0)$, если между началом и этой точкой он оси не касается. Сколько путей Дика в $(2n,0)$ впервые возвращаются на~ось в~точке~$(2k,0)$?}
\punkt{б}{}{Выразите $C_{n+1}$ через $C_0,\ldots,C_n$.}
\punkt{в}{}{Найдите ответ в задаче 1 при $a=n+1$, $b=n$.}
\end{samepage}\par\vspace{9pt}

\tekst{Дадим ещё три определения.}

\celoe{\tekst{$\bullet$~\term{Триангуляция} выпуклого многоугольника~— разрезание его непересекающимися диагоналями на треугольники. Множество триангуляций правильного $(n+2)$-угольника обозначим $\mathcal T_n$; триангуляции, переходящие друг в друга при повороте или симметрии, различаются. На рисунке~— все элементы $\mathcal T_3$.}\risunok{\begin{tikzpicture}[scale=1.1] \draw[line width=1.1pt,line join=round] (0.000,1.100) -- (-1.046,0.340) -- (-0.647,-0.890) -- (0.647,-0.890) -- (1.046,0.340) -- cycle; \draw[line width=0.8pt] (0.000,1.100) -- (-0.647,-0.890) (0.000,1.100) -- (0.647,-0.890); \fill (0.000,1.100) circle (1.7pt); \fill (-1.046,0.340) circle (1.7pt); \fill (-0.647,-0.890) circle (1.7pt); \fill (0.647,-0.890) circle (1.7pt); \fill (1.046,0.340) circle (1.7pt); \draw[line width=1.1pt,line join=round] (3.500,1.100) -- (2.454,0.340) -- (2.853,-0.890) -- (4.147,-0.890) -- (4.546,0.340) -- cycle; \draw[line width=0.8pt] (2.454,0.340) -- (4.147,-0.890) (2.454,0.340) -- (4.546,0.340); \fill (3.500,1.100) circle (1.7pt); \fill (2.454,0.340) circle (1.7pt); \fill (2.853,-0.890) circle (1.7pt); \fill (4.147,-0.890) circle (1.7pt); \fill (4.546,0.340) circle (1.7pt); \draw[line width=1.1pt,line join=round] (7.000,1.100) -- (5.954,0.340) -- (6.353,-0.890) -- (7.647,-0.890) -- (8.046,0.340) -- cycle; \draw[line width=0.8pt] (6.353,-0.890) -- (8.046,0.340) (6.353,-0.890) -- (7.000,1.100); \fill (7.000,1.100) circle (1.7pt); \fill (5.954,0.340) circle (1.7pt); \fill (6.353,-0.890) circle (1.7pt); \fill (7.647,-0.890) circle (1.7pt); \fill (8.046,0.340) circle (1.7pt); \draw[line width=1.1pt,line join=round] (10.500,1.100) -- (9.454,0.340) -- (9.853,-0.890) -- (11.147,-0.890) -- (11.546,0.340) -- cycle; \draw[line width=0.8pt] (11.147,-0.890) -- (10.500,1.100) (11.147,-0.890) -- (9.454,0.340); \fill (10.500,1.100) circle (1.7pt); \fill (9.454,0.340) circle (1.7pt); \fill (9.853,-0.890) circle (1.7pt); \fill (11.147,-0.890) circle (1.7pt); \fill (11.546,0.340) circle (1.7pt); \draw[line width=1.1pt,line join=round] (14.000,1.100) -- (12.954,0.340) -- (13.353,-0.890) -- (14.647,-0.890) -- (15.046,0.340) -- cycle; \draw[line width=0.8pt] (15.046,0.340) -- (12.954,0.340) (15.046,0.340) -- (13.353,-0.890); \fill (14.000,1.100) circle (1.7pt); \fill (12.954,0.340) circle (1.7pt); \fill (13.353,-0.890) circle (1.7pt); \fill (14.647,-0.890) circle (1.7pt); \fill (15.046,0.340) circle (1.7pt); \end{tikzpicture}}}

\celoe{\tekst{$\bullet$~\term{Двоичное дерево}~— дерево с выделенной вершиной (корнем), в котором у каждой вершины либо нет потомков, либо их ровно два, причём один из них назван левым, а другой~— правым. Вершины с потомками~— \term{развилки} (крупные точки). Множество двоичных деревьев с $n$ развилками обозначим $\mathcal D_n$. На рисунке~— все элементы $\mathcal D_3$.}\risunok{\begin{tikzpicture}[scale=1.1] \draw[line width=0.9pt] (0.000,0.000) -- (-0.800,-0.620); \draw[line width=0.9pt] (0.000,0.000) -- (0.800,-0.620); \draw[line width=0.9pt] (0.800,-0.620) -- (0.360,-1.240); \draw[line width=0.9pt] (0.800,-0.620) -- (1.240,-1.240); \draw[line width=0.9pt] (1.240,-1.240) -- (0.998,-1.860); \draw[line width=0.9pt] (1.240,-1.240) -- (1.482,-1.860); \fill (0.000,0.000) circle (2.3pt); \fill (-0.800,-0.620) circle (1.6pt); \fill (0.800,-0.620) circle (2.3pt); \fill (0.360,-1.240) circle (1.6pt); \fill (1.240,-1.240) circle (2.3pt); \fill (0.998,-1.860) circle (1.6pt); \fill (1.482,-1.860) circle (1.6pt); \draw[line width=0.9pt] (3.500,0.000) -- (2.700,-0.620); \draw[line width=0.9pt] (3.500,0.000) -- (4.300,-0.620); \draw[line width=0.9pt] (4.300,-0.620) -- (3.860,-1.240); \draw[line width=0.9pt] (3.860,-1.240) -- (3.618,-1.860); \draw[line width=0.9pt] (3.860,-1.240) -- (4.102,-1.860); \draw[line width=0.9pt] (4.300,-0.620) -- (4.740,-1.240); \fill (3.500,0.000) circle (2.3pt); \fill (2.700,-0.620) circle (1.6pt); \fill (4.300,-0.620) circle (2.3pt); \fill (3.860,-1.240) circle (2.3pt); \fill (3.618,-1.860) circle (1.6pt); \fill (4.102,-1.860) circle (1.6pt); \fill (4.740,-1.240) circle (1.6pt); \draw[line width=0.9pt] (7.000,0.000) -- (6.200,-0.620); \draw[line width=0.9pt] (6.200,-0.620) -- (5.760,-1.240); \draw[line width=0.9pt] (6.200,-0.620) -- (6.640,-1.240); \draw[line width=0.9pt] (7.000,0.000) -- (7.800,-0.620); \draw[line width=0.9pt] (7.800,-0.620) -- (7.360,-1.240); \draw[line width=0.9pt] (7.800,-0.620) -- (8.240,-1.240); \fill (7.000,0.000) circle (2.3pt); \fill (6.200,-0.620) circle (2.3pt); \fill (5.760,-1.240) circle (1.6pt); \fill (6.640,-1.240) circle (1.6pt); \fill (7.800,-0.620) circle (2.3pt); \fill (7.360,-1.240) circle (1.6pt); \fill (8.240,-1.240) circle (1.6pt); \draw[line width=0.9pt] (10.500,0.000) -- (9.700,-0.620); \draw[line width=0.9pt] (9.700,-0.620) -- (9.260,-1.240); \draw[line width=0.9pt] (9.700,-0.620) -- (10.140,-1.240); \draw[line width=0.9pt] (10.140,-1.240) -- (9.898,-1.860); \draw[line width=0.9pt] (10.140,-1.240) -- (10.382,-1.860); \draw[line width=0.9pt] (10.500,0.000) -- (11.300,-0.620); \fill (10.500,0.000) circle (2.3pt); \fill (9.700,-0.620) circle (2.3pt); \fill (9.260,-1.240) circle (1.6pt); \fill (10.140,-1.240) circle (2.3pt); \fill (9.898,-1.860) circle (1.6pt); \fill (10.382,-1.860) circle (1.6pt); \fill (11.300,-0.620) circle (1.6pt); \draw[line width=0.9pt] (14.000,0.000) -- (13.200,-0.620); \draw[line width=0.9pt] (13.200,-0.620) -- (12.760,-1.240); \draw[line width=0.9pt] (12.760,-1.240) -- (12.518,-1.860); \draw[line width=0.9pt] (12.760,-1.240) -- (13.002,-1.860); \draw[line width=0.9pt] (13.200,-0.620) -- (13.640,-1.240); \draw[line width=0.9pt] (14.000,0.000) -- (14.800,-0.620); \fill (14.000,0.000) circle (2.3pt); \fill (13.200,-0.620) circle (2.3pt); \fill (12.760,-1.240) circle (2.3pt); \fill (12.518,-1.860) circle (1.6pt); \fill (13.002,-1.860) circle (1.6pt); \fill (13.640,-1.240) circle (1.6pt); \fill (14.800,-0.620) circle (1.6pt); \end{tikzpicture}}}

\sboku{8.2cm}{\begin{tikzpicture}[>=stealth,scale=0.93] \draw[line width=2.4pt,gray!45,line cap=round] (0,0) -- (3.15,0) (4.25,0) -- (8.4,0); \draw[line width=1.4pt,gray!60] (3.15,0) -- (3.15,-2.75) -- (4.25,-2.75) -- (4.25,0); \fill[gray!12] (3.17,-0.02) rectangle (4.23,-2.73); \draw[line width=1.4pt,gray!60] (3.15,0) -- (3.15,-2.75) -- (4.25,-2.75) -- (4.25,0); \draw[line width=0.8pt,rounded corners=1.5pt,fill=white] (1.49,0.05) rectangle (2.31,0.55); \node[font=\small] at (1.90,0.30) {2}; \draw[line width=0.8pt,rounded corners=1.5pt,fill=white] (3.29,-2.68) rectangle (4.11,-2.18); \node[font=\small] at (3.70,-2.43) {1}; \draw[line width=0.8pt,rounded corners=1.5pt,fill=white] (3.29,-2.10) rectangle (4.11,-1.60); \node[font=\small] at (3.70,-1.85) {3}; \draw[line width=0.8pt,rounded corners=1.5pt,fill=white] (5.19,0.05) rectangle (6.01,0.55); \node[font=\small] at (5.60,0.30) {4}; \draw[->,line width=0.8pt] (6.4,0.95) .. controls (4.6,1.1) and (3.95,0.6) .. (3.95,-0.45); \draw[->,line width=0.8pt] (3.45,-0.45) .. controls (3.45,0.6) and (2.9,1.1) .. (1.2,0.95); \node[font=\small\itshape,gray] at (8.0,0.45) {вход}; \node[font=\small\itshape,gray] at (0.45,0.45) {выход}; \node[font=\small\itshape,gray] at (5.0,-2.4) {тупик}; \end{tikzpicture}}{\tekst{$\bullet$~Рассмотрим железнодорожный тупик, справа от которого стоят вагоны, пронумерованные числами от 1 до $n$. Вагоны заезжают в тупик по одному в порядке номеров; в любой момент вагон, стоящий в тупике ближе всех к выезду, может выехать налево. Строка номеров вагонов в порядке выезда называется \term{стековой перестановкой}. Множество стековых перестановок длины $n$ обозначим $\mathcal S_n$.\\
Например, на рисунке $n=4$: вагон 2 уже выехал, 1 и 3 в тупике; если выедут 3 и 1, а затем 4, получится $2314$.}}

\begin{samepage}
\blok{\nomer{5}{\dagger}\hspace{\zazor}Пусть $\mathcal P_n$~— множество путей Дика в $(2n,0)$. Для каждого из множеств $\mathcal P_n$, $\mathcal D_n$, $\mathcal T_n$, $\mathcal S_n$:}
\punkt{а}{}{выпишите (нарисуйте) все его элементы при $n\leqslant 4$;}
\punkt{б}{}{выразите число его элементов при $n+1$ через числа элементов при меньших $n$;}
\punkt{в}{}{постройте биекции между ним и каждым из остальных.}
\end{samepage}\par\vspace{9pt}

\zad{6}{}{При каких $n$ число $C_n$ нечётно?}

\begin{samepage}
\blok{\nomer{7}{}\hspace{\zazor}«Бесконечный многочлен» $C(x)=C_0+C_1x+C_2x^2+\ldots$ назовём \term{производящей функцией} чисел~Каталана.}
\punkt{а}{}{Докажите, что $x\,C(x)^2=C(x)-1$.}
\punkt{б}{}{Найдите $C(x)$.}
\end{samepage}\par\vspace{9pt}

\newpage
\zvezdy

\tekst{Дальше~— несколько разных доказательств предположения из задачи 3в.}

\sboku{7.9cm}{\begin{tikzpicture}[scale=0.29] \draw[step=1,gray!28,line width=0.3pt] (0,-3) grid (22,3); \draw[gray!70,line width=0.6pt,->] (0,0) -- (22.4,0); \draw[line width=1.3pt,line join=round,line cap=round] (0,0) -- (1,1) -- (2,0) -- (3,-1) -- (4,0) -- (5,1) -- (6,2) -- (7,1) -- (8,0) -- (9,-1) -- (10,-2) -- (11,-1) -- (12,0) -- (13,1) -- (14,2) -- (15,1) -- (16,0) -- (17,-1) -- (18,0) -- (19,1) -- (20,2) -- (21,1) -- (22,0); \draw[dashed,line width=0.9pt] (0,-1) -- (22,-1) node[right,font=\small]{$y=-1$}; \fill (0,0) circle (3.5pt); \fill (22,0) circle (3.5pt); \end{tikzpicture}}{\begin{samepage}
\blok{\nomer{8}{\dagger}\hspace{\zazor}Путь из $(0,0)$ назовём \term{плохим}, если он хотя бы раз опускается ниже оси (на рисунке~— плохой путь и прямая $y=-1$).}
\punkt{а}{}{Сколько плохих путей из $(0,0)$ в $(4,0)$? в $(6,0)$?}
\punkt{б}{}{Найдите $C_n$.}
\punkt{в}{}{Сколько путей Дика ведут в точку $(m,h)$?}
\punkt{г}{}{Найдите ответ в задаче 1.}
\end{samepage}\par}

\zad{9}{\circ}{С какой вероятностью на выборах из задачи 1 А ни разу не отставал от Б, если $a\geqslant b$?}

\sboku{3.9cm}{\begin{tikzpicture}[>=stealth] \draw[gray!60,line width=0.7pt] (0,0) circle (1.15); \fill (0.499,1.036) circle (1.6pt); \node[font=\normalsize] at (0.673,1.397) {$+1$}; \fill (1.121,0.256) circle (1.6pt); \node[font=\normalsize] at (1.511,0.345) {$-1$}; \fill (0.899,-0.717) circle (1.6pt); \node[font=\normalsize] at (1.212,-0.966) {$+1$}; \fill (0.000,-1.150) circle (1.6pt); \node[font=\normalsize] at (0.000,-1.550) {$+1$}; \fill (-0.899,-0.717) circle (1.6pt); \node[font=\normalsize] at (-1.212,-0.966) {$-1$}; \fill (-1.121,0.256) circle (1.6pt); \node[font=\normalsize] at (-1.511,0.345) {$+1$}; \fill (-0.499,1.036) circle (1.6pt); \node[font=\normalsize] at (-0.673,1.397) {$-1$}; \draw[line width=0.8pt,fill=white] (90:1.15) circle (2.6pt); \draw[->,line width=0.8pt] (84:0.82) arc (84:-46.6:0.82); \node[font=\normalsize] at (0,0) {$1$}; \end{tikzpicture}}{\begin{samepage}
\blok{\nomer{10}{}\hspace{\zazor}(\textit{Циклическая лемма.}) По кругу стоят $a$ чисел $+1$ и $b$ чисел $-1$, $a>b$. Начнём с одного из мест и по часовой стрелке прибавляем каждое число к сумме предыдущих.}
\punkt{а}{}{Сколько существует начальных мест, для которых все суммы положительны?}
\punkt{б}{}{Выведите отсюда формулу для $C_n$ другим способом.}
\end{samepage}\par}

\celoe{\begin{samepage}
\blok{\nomer{11}{}\hspace{\zazor}Звено пути назовём \term{верхним}, если хотя бы один его конец выше оси.}
\punkt{а}{}{Сколько путей из $(0,0)$ в $(6,0)$ имеют $0,2,4,6$ верхних звеньев?}
\punkt{б}{}{Сколько путей из $(0,0)$ в $(2n,0)$ имеют ровно $2k$ верхних звеньев?}
\punkt{в}{}{Выведите отсюда формулу для $C_n$.}
\end{samepage}\risunok{\begin{tikzpicture}[scale=0.5] \draw[step=1,gray!28,line width=0.3pt] (0,-2) grid (4,2); \draw[gray!70,line width=0.6pt] (0,0) -- (4,0); \draw[line width=2.6pt,line cap=round] (0,0) -- (1,1); \draw[line width=2.6pt,line cap=round] (1,1) -- (2,2); \draw[line width=2.6pt,line cap=round] (2,2) -- (3,1); \draw[line width=2.6pt,line cap=round] (3,1) -- (4,0); \draw[step=1,gray!28,line width=0.3pt] (5.3,-2) grid (9.3,2); \draw[gray!70,line width=0.6pt] (5.3,0) -- (9.3,0); \draw[line width=2.6pt,line cap=round] (5.3,0) -- (6.3,1); \draw[line width=2.6pt,line cap=round] (6.3,1) -- (7.3,0); \draw[line width=2.6pt,line cap=round] (7.3,0) -- (8.3,1); \draw[line width=2.6pt,line cap=round] (8.3,1) -- (9.3,0); \draw[step=1,gray!28,line width=0.3pt] (10.6,-2) grid (14.6,2); \draw[gray!70,line width=0.6pt] (10.6,0) -- (14.6,0); \draw[line width=2.6pt,line cap=round] (10.6,0) -- (11.6,1); \draw[line width=2.6pt,line cap=round] (11.6,1) -- (12.6,0); \draw[line width=0.9pt,gray] (12.6,0) -- (13.6,-1); \draw[line width=0.9pt,gray] (13.6,-1) -- (14.6,0); \draw[step=1,gray!28,line width=0.3pt] (15.9,-2) grid (19.9,2); \draw[gray!70,line width=0.6pt] (15.9,0) -- (19.9,0); \draw[line width=0.9pt,gray] (15.9,0) -- (16.9,-1); \draw[line width=0.9pt,gray] (16.9,-1) -- (17.9,0); \draw[line width=2.6pt,line cap=round] (17.9,0) -- (18.9,1); \draw[line width=2.6pt,line cap=round] (18.9,1) -- (19.9,0); \draw[step=1,gray!28,line width=0.3pt] (21.2,-2) grid (25.2,2); \draw[gray!70,line width=0.6pt] (21.2,0) -- (25.2,0); \draw[line width=0.9pt,gray] (21.2,0) -- (22.2,-1); \draw[line width=0.9pt,gray] (22.2,-1) -- (23.2,0); \draw[line width=0.9pt,gray] (23.2,0) -- (24.2,-1); \draw[line width=0.9pt,gray] (24.2,-1) -- (25.2,0); \draw[step=1,gray!28,line width=0.3pt] (26.5,-2) grid (30.5,2); \draw[gray!70,line width=0.6pt] (26.5,0) -- (30.5,0); \draw[line width=0.9pt,gray] (26.5,0) -- (27.5,-1); \draw[line width=0.9pt,gray] (27.5,-1) -- (28.5,-2); \draw[line width=0.9pt,gray] (28.5,-2) -- (29.5,-1); \draw[line width=0.9pt,gray] (29.5,-1) -- (30.5,0); \end{tikzpicture}}}

\sboku{3.6cm}{\begin{tikzpicture} \draw[line width=1.1pt,line join=round] (-0.542,-1.126) -- (0.542,-1.126) -- (1.219,-0.278) -- (0.977,0.779) -- (0.000,1.250) -- (-0.977,0.779) -- (-1.219,-0.278) -- cycle; \draw[line width=0.8pt] (0.542,-1.126) -- (0.977,0.779); \draw[line width=0.8pt] (0.977,0.779) -- (-1.219,-0.278); \draw[line width=0.8pt] (-1.219,-0.278) -- (0.542,-1.126); \draw[line width=0.8pt] (0.977,0.779) -- (-0.977,0.779); \draw[line width=3pt,line cap=round] (-0.542,-1.126) -- (0.542,-1.126); \draw[line width=3pt,line cap=round,gray!55] (-0.977,0.779) -- (-1.219,-0.278); \fill (-0.542,-1.126) circle (1.7pt); \fill (0.542,-1.126) circle (1.7pt); \fill (1.219,-0.278) circle (1.7pt); \fill (0.977,0.779) circle (1.7pt); \fill (0.000,1.250) circle (1.7pt); \fill (-0.977,0.779) circle (1.7pt); \fill (-1.219,-0.278) circle (1.7pt); \end{tikzpicture}}{\zad{12}{}{Постройте соответствие, которое каждому элементу $\mathcal T_{n+1}$ с отмеченной стороной, отличной от нижней, сопоставляет элемент $\mathcal T_n$ так, что каждый получается ровно $4n+2$ раза. Выведите отсюда формулу для~$C_n$ ещё одним~способом.}\par}

\zvezdy

\begin{samepage}
\blok{\nomer{13}{\circ}\hspace{\zazor}Монету бросают $2n$ раз. \term{Момент} $k$~— сразу после $k$-го броска, момент $0$~— начало. С какой вероятностью орлов и решек ни в какой момент, кроме $0$, не будет поровну?}
\punkt{а}{}{Найдите ответ при $2n=2,4,6,8$.}
\punkt{б}{}{Докажите, что таких исходов вдвое больше, чем путей Дика длины $2n-1$ (с любым концом).}
\punkt{в}{}{Найдите ответ при любом $n$.}
\end{samepage}\par\vspace{9pt}

\celoe{\begin{samepage}
\zadp{14}{\dagger}
\punkt{а}{}{Сколько путей из $(0,0)$ в $(6,0)$ имеют самую низкую точку на высоте $0,-1,-2,-3$?}
\punkt{б}{}{Постройте биекцию между путями длины $2n$, которые кончаются на оси, и путями Дика длины $2n$. Выведите отсюда ответ в задаче 13 другим способом.}
\end{samepage}\risunok{\begin{tikzpicture}[scale=0.36] \draw[step=1,gray!28,line width=0.3pt] (0,-4) grid (18,3); \draw[gray!70,line width=0.6pt,->] (0,0) -- (18.4,0); \draw[line width=1.3pt,line join=round,line cap=round] (0,0) -- (1,1) -- (2,0) -- (3,-1) -- (4,0) -- (5,-1) -- (6,-2) -- (7,-3) -- (8,-2) -- (9,-1) -- (10,-2) -- (11,-1) -- (12,0) -- (13,1) -- (14,2) -- (15,1) -- (16,2) -- (17,1) -- (18,0); \fill (7,-3) circle (4pt); \end{tikzpicture}}}

\begin{samepage}
\zadp{15}{}
\punkt{а}{}{Монету бросают 10\,000 раз. Докажите, что с вероятностью больше 0,99 орлов и решек после начала хотя бы раз будет поровну.}
\punkt{б}{}{Докажите, что ответ $u_n$ в задаче 13 удовлетворяет неравенствам \[\frac{1}{2\sqrt n}\leqslant u_n\leqslant\frac{1}{\sqrt{2n+1}}.\]}
\punkt{в}{}{С какой вероятностью орлов и решек впервые после начала станет поровну в момент $2n$? Докажите, что \[\frac{C_0}{4}+\frac{C_1}{4^2}+\ldots+\frac{C_{N-1}}{4^N}=\frac{1-u_N}{2}.\]}
\punkt{г}{}{Найдите $C(\frac14)$, где $C(x)$~— производящая функция из задачи 7. Как это согласуется с~пунктом~в)?}
\end{samepage}\par\vspace{9pt}

\zvezdy

\celoe{\begin{samepage}
\blok{\nomer{16}{\circ}\hspace{\zazor}Двое играют в орлянку: монету бросают несколько раз, орёл~— очко первому, решка~— второму; \term{путь игры} идёт звеном $\nearrow$ при орле и $\searrow$ при решке.}
\punkt{а}{}{При 4 и 6 бросках найдите вероятность того, что последний равный счёт был в~момент~$0,2,4,\ldots$}
\punkt{б}{}{Найдите вероятность того, что при $2n$ бросках последний равный счёт был в момент $2k$.}
\punkt{в}{}{Игра длится 100 бросков. Какова вероятность того, что последний момент равного счёта окажется среди моментов от 0 до 10? от 90 до 100? от 45 до 55?}
\end{samepage}\risunok{\begin{tikzpicture}[scale=0.34] \draw[step=1,gray!28,line width=0.3pt] (0,-6) grid (40,1); \draw[gray!70,line width=0.6pt,->] (0,0) -- (40.4,0); \draw[line width=1.3pt,line join=round,line cap=round] (0,0) -- (1,-1) -- (2,0) -- (3,-1) -- (4,0) -- (5,-1) -- (6,0) -- (7,-1) -- (8,-2) -- (9,-3) -- (10,-4) -- (11,-3) -- (12,-4) -- (13,-3) -- (14,-4) -- (15,-3) -- (16,-4) -- (17,-5) -- (18,-4) -- (19,-3) -- (20,-4) -- (21,-3) -- (22,-2) -- (23,-1) -- (24,-2) -- (25,-1) -- (26,-2) -- (27,-1) -- (28,-2) -- (29,-1) -- (30,-2) -- (31,-3) -- (32,-4) -- (33,-5) -- (34,-4) -- (35,-3) -- (36,-4) -- (37,-3) -- (38,-4) -- (39,-3) -- (40,-4); \draw[line width=1.1pt] (6,0) circle (0.6); \end{tikzpicture}}}

\celoe{\sboku{7.6cm}{\begin{tikzpicture}[>=stealth] \draw[gray!28,line width=0.3pt] (0,0.72) -- (6.2,0.72); \draw[gray!28,line width=0.3pt] (0,1.44) -- (6.2,1.44); \draw[gray!28,line width=0.3pt] (0,2.16) -- (6.2,2.16); \draw[gray!28,line width=0.3pt] (0,2.88) -- (6.2,2.88); \draw[gray!28,line width=0.3pt] (1.55,0) -- (1.55,2.88); \draw[gray!28,line width=0.3pt] (3.1,0) -- (3.1,2.88); \draw[gray!28,line width=0.3pt] (4.65,0) -- (4.65,2.88); \draw[gray!70,line width=0.6pt,->] (0,0) -- (6.55,0) node[right,font=\small]{$x$}; \draw[gray!70,line width=0.6pt,->] (0,0) -- (0,3.23) node[above,font=\small]{$y$}; \draw[dashed,gray,line width=0.5pt] (6.2,0) -- (6.2,2.88); \draw[line width=1.3pt,line join=round] (0.0409,2.8304) -- (0.0430,2.7611) -- (0.0493,2.5804) -- (0.0598,2.3459) -- (0.0743,2.1056) -- (0.0930,1.8851) -- (0.1158,1.6930) -- (0.1425,1.5293) -- (0.1732,1.3909) -- (0.2076,1.2739) -- (0.2458,1.1745) -- (0.2877,1.0895) -- (0.3330,1.0165) -- (0.3818,0.9534) -- (0.4338,0.8984) -- (0.4889,0.8504) -- (0.5470,0.8081) -- (0.6079,0.7707) -- (0.6714,0.7375) -- (0.7374,0.7080) -- (0.8057,0.6816) -- (0.8761,0.6579) -- (0.9483,0.6367) -- (1.0223,0.6176) -- (1.0978,0.6004) -- (1.1746,0.5849) -- (1.2525,0.5708) -- (1.3312,0.5581) -- (1.4106,0.5467) -- (1.4904,0.5363) -- (1.5705,0.5270) -- (1.6505,0.5185) -- (1.7303,0.5109) -- (1.8097,0.5041) -- (1.8885,0.4980) -- (1.9663,0.4925) -- (2.0431,0.4876) -- (2.1186,0.4832) -- (2.1926,0.4794) -- (2.2649,0.4760) -- (2.3352,0.4730) -- (2.4035,0.4704) -- (2.4695,0.4682) -- (2.5330,0.4662) -- (2.5939,0.4646) -- (2.6520,0.4632) -- (2.7071,0.4621) -- (2.7591,0.4612) -- (2.8079,0.4604) -- (2.8532,0.4598) -- (2.8951,0.4594) -- (2.9333,0.4590) -- (2.9678,0.4588) -- (2.9984,0.4586) -- (3.0251,0.4585) -- (3.0479,0.4584) -- (3.0666,0.4584) -- (3.0812,0.4584) -- (3.0916,0.4584) -- (3.0979,0.4584) -- (3.1000,0.4584) -- (3.1021,0.4584) -- (3.1084,0.4584) -- (3.1188,0.4584) -- (3.1334,0.4584) -- (3.1521,0.4584) -- (3.1749,0.4585) -- (3.2016,0.4586) -- (3.2322,0.4588) -- (3.2667,0.4590) -- (3.3049,0.4594) -- (3.3468,0.4598) -- (3.3921,0.4604) -- (3.4409,0.4612) -- (3.4929,0.4621) -- (3.5480,0.4632) -- (3.6061,0.4646) -- (3.6670,0.4662) -- (3.7305,0.4682) -- (3.7965,0.4704) -- (3.8648,0.4730) -- (3.9351,0.4760) -- (4.0074,0.4794) -- (4.0814,0.4832) -- (4.1569,0.4876) -- (4.2337,0.4925) -- (4.3115,0.4980) -- (4.3903,0.5041) -- (4.4697,0.5109) -- (4.5495,0.5185) -- (4.6295,0.5270) -- (4.7096,0.5363) -- (4.7894,0.5467) -- (4.8688,0.5581) -- (4.9475,0.5708) -- (5.0254,0.5849) -- (5.1022,0.6004) -- (5.1777,0.6176) -- (5.2517,0.6367) -- (5.3239,0.6579) -- (5.3943,0.6816) -- (5.4626,0.7080) -- (5.5286,0.7375) -- (5.5921,0.7707) -- (5.6530,0.8081) -- (5.7111,0.8504) -- (5.7662,0.8984) -- (5.8182,0.9534) -- (5.8670,1.0165) -- (5.9123,1.0895) -- (5.9542,1.1745) -- (5.9924,1.2739) -- (6.0268,1.3909) -- (6.0575,1.5293) -- (6.0842,1.6930) -- (6.1070,1.8851) -- (6.1257,2.1056) -- (6.1402,2.3459) -- (6.1507,2.5804) -- (6.1570,2.7611) -- (6.1591,2.8304); \node[below,font=\small] at (0,0) {$0$}; \node[below,font=\small] at (3.1,0) {$\frac12$}; \node[below,font=\small] at (6.2,0) {$1$}; \node[left,font=\small] at (0,0.72) {$1$}; \node[left,font=\small] at (0,1.44) {$2$}; \node[left,font=\small] at (0,2.16) {$3$}; \node[left,font=\small] at (0,2.88) {$4$}; \node[font=\small,fill=white,inner sep=1.5pt] at (3.1,1.692) {$y=\dfrac{1}{\pi\sqrt{x(1-x)}}$}; \end{tikzpicture}}{\begin{samepage}
\blok{\nomer{17}{\dagger}\hspace{\zazor}Игра длится $2n$ бросков. Будем говорить, что первый игрок \term{ведёт} во время броска, если соответствующее звено пути игры верхнее (см. задачу 11).}
\punkt{а}{}{Сколько путей длины 4, 6, 8 имеют $2k$ верхних звеньев?}
\punkt{б}{}{Постройте биекцию между путями длины $2n$ с $2k$ верхними звеньями и путями длины $2n$, у которых в последний раз счёт был равным в момент $2k$.}
\punkt{в}{}{Смоделируйте на компьютере 10\,000 игр по 1000 бросков. Постройте гистограмму доли времени, когда вёл первый, и нанесите на тот же чертёж график функции справа. В какой доле игр первый вёл больше 90\% времени?}
\end{samepage}\par}\risunok{\begin{tikzpicture}[scale=0.38] \draw[step=1,gray!28,line width=0.3pt] (0,-3) grid (24,4); \draw[gray!70,line width=0.6pt,->] (0,0) -- (24.4,0); \draw[line width=2.6pt,line cap=round] (0,0) -- (1,1); \draw[line width=2.6pt,line cap=round] (1,1) -- (2,2); \draw[line width=2.6pt,line cap=round] (2,2) -- (3,1); \draw[line width=2.6pt,line cap=round] (3,1) -- (4,0); \draw[line width=0.9pt,gray] (4,0) -- (5,-1); \draw[line width=0.9pt,gray] (5,-1) -- (6,0); \draw[line width=0.9pt,gray] (6,0) -- (7,-1); \draw[line width=0.9pt,gray] (7,-1) -- (8,-2); \draw[line width=0.9pt,gray] (8,-2) -- (9,-1); \draw[line width=0.9pt,gray] (9,-1) -- (10,0); \draw[line width=2.6pt,line cap=round] (10,0) -- (11,1); \draw[line width=2.6pt,line cap=round] (11,1) -- (12,2); \draw[line width=2.6pt,line cap=round] (12,2) -- (13,1); \draw[line width=2.6pt,line cap=round] (13,1) -- (14,2); \draw[line width=2.6pt,line cap=round] (14,2) -- (15,1); \draw[line width=2.6pt,line cap=round] (15,1) -- (16,0); \draw[line width=0.9pt,gray] (16,0) -- (17,-1); \draw[line width=0.9pt,gray] (17,-1) -- (18,-2); \draw[line width=0.9pt,gray] (18,-2) -- (19,-1); \draw[line width=0.9pt,gray] (19,-1) -- (20,0); \draw[line width=2.6pt,line cap=round] (20,0) -- (21,1); \draw[line width=2.6pt,line cap=round] (21,1) -- (22,2); \draw[line width=2.6pt,line cap=round] (22,2) -- (23,3); \draw[line width=2.6pt,line cap=round] (23,3) -- (24,2); \end{tikzpicture}}}

\celoe{\begin{samepage}
\blok{\nomer{18}{}\hspace{\zazor}\term{Сменой лидера} назовём момент, когда счёт равный, причём перед ним путь игры шёл по одну сторону от оси, а после~— по другую (на рисунке смены отмечены крестиками, равный счёт без смены~— кружками). Игра длится $2n$ бросков.}
\punkt{а}{}{Сколько путей длины 4, 6, 8 имеют $0,1,2,\ldots$ смен лидера?}
\punkt{б}{}{Докажите, что вероятность ровно $r$ смен равна $2\binom{2n-1}{n-1-r}\big/2^{2n-1}$. Какое число смен наиболее вероятно?}
\punkt{в}{}{Во сколько раз «лидер ни разу не сменился» вероятнее, чем «после начала счёт ни разу не~был~равным»?}
\end{samepage}\risunok{\begin{tikzpicture}[scale=0.38] \draw[step=1,gray!28,line width=0.3pt] (0,-3) grid (28,3); \draw[gray!70,line width=0.6pt,->] (0,0) -- (28.4,0); \draw[line width=1.3pt,line join=round,line cap=round] (0,0) -- (1,1) -- (2,0) -- (3,1) -- (4,2) -- (5,1) -- (6,0) -- (7,-1) -- (8,0) -- (9,-1) -- (10,-2) -- (11,-1) -- (12,0) -- (13,1) -- (14,2) -- (15,1) -- (16,0) -- (17,1) -- (18,2) -- (19,1) -- (20,0) -- (21,-1) -- (22,-2) -- (23,-1) -- (24,0) -- (25,-1) -- (26,0) -- (27,1) -- (28,2); \draw[line width=1.1pt] (2,0) circle (0.45); \draw[line width=1.5pt] (5.55,-0.45) -- (6.45,0.45) (5.55,0.45) -- (6.45,-0.45); \draw[line width=1.1pt] (8,0) circle (0.45); \draw[line width=1.5pt] (11.55,-0.45) -- (12.45,0.45) (11.55,0.45) -- (12.45,-0.45); \draw[line width=1.1pt] (16,0) circle (0.45); \draw[line width=1.5pt] (19.55,-0.45) -- (20.45,0.45) (19.55,0.45) -- (20.45,-0.45); \draw[line width=1.1pt] (24,0) circle (0.45); \draw[line width=1.5pt] (25.55,-0.45) -- (26.45,0.45) (25.55,0.45) -- (26.45,-0.45); \end{tikzpicture}}}

\end{document}
